\documentclass[10pt,a4paper]{article}
\usepackage{amsmath,amssymb,amsthm,graphicx}
\usepackage{xfrac}
\newenvironment{enumeratei}
{\bgroup\def\theenumi{\roman{enumi}}\def\theenumii{\arabic{enumii}}\begin{enumerate}}
{\end{enumerate}\egroup}

\hyphenation{addi-tion-al belong belongs denote exchan-ge exchan-ged exchan-ges}

\newcommand{\psfrac}[2]{\sfrac{(#1)}{#2}}

\newcommand\mto{\mathchoice{\longmapsto}{\mapsto}{\mapsto}{\mapsto}}
\let\oldbigsqcup\bigsqcup
\renewcommand{\bigsqcup}{\mathop{\textstyle\oldbigsqcup}\displaylimits}

\let\epsilon\varepsilon

\DeclareMathOperator{\Ess}{Ess}
\DeclareMathOperator{\Leb}{Leb}

\usepackage{tikz}
\usetikzlibrary{plotmarks}

\theoremstyle{plain}
\newtheorem{theorem}{Theorem}[section]
\newtheorem{lemma}[theorem]{Lemma}
\newtheorem{proposition}[theorem]{Proposition}
\newtheorem{corollary}[theorem]{Corollary}
\newtheorem{claim}{Claim}

\theoremstyle{definition}
\newtheorem{remark}[theorem]{Remark}
\newtheorem{example}[theorem]{Example}

\newcommand{\A}{\mathcal{A}}
\newcommand{\B}{\mathcal{B}}
\newcommand{\I}{\mathcal{I}}

\newcommand{\Z}{\mathbb{Z}}
\newcommand{\N}{\mathbb{N}}
\newcommand{\R}{\mathbb{R}}
\title{An ergodic symmetric interval exchange transformation has a centre outside every non-trivial connection}
\author{Przemyslaw Berk, Frank Trujillo}
\date{doi:10.5802/jep.302; CC BY4.0}
\begin{document}
\maketitle
\noindent\textit{Editorial note.} The counted statement is the author's labelled lemma environment at source0.tex lines 728-731; its proof environment opens at line 733, the line after the blank separator, so the pairing is by adjacency and the retained proof range 733-743 is the whole environment.\par\medskip
\noindent\textit{Editorial note.} Three statements are retained as context because the counted proof cites them: cor: existenceofinfiniteorbit (source lines 220-222), cor:centers\_connections (source lines 686-689) and lem: sym\_endpoints (source lines 693-704).\par\medskip
\noindent\textit{Editorial note.} This source is new to the material-authoring work and its preamble is built from the author's own macro and theorem block (source lines 8-45) rather than from a prepared adapter.\par\medskip
\noindent\textit{Editorial note.} This source has no label-number dump and the retained counter values are left empty, so the printed numbering of the card is the wrapper's default; the identity of the counted item is its author label.\par\medskip
\noindent\textit{Editorial scope.} One original result of the pinned publisher article, counted once: existence of a centre lying in no non-trivial connection for an ergodic symmetric interval exchange transformation (author label lem:sym\_int\_disjoint).\par\medskip
\noindent\textit{Allowed standard prerequisites (editorial).}\begin{itemize}
\item Connection. A connection of an interval exchange transformation is a finite orbit fragment joining two discontinuities.. The statement asserts the existence of a centre lying in no non-trivial connection.
\item Centres and connections. Any non-trivial irreducible connection contains at most one point from the set of middle points of the exchanged intervals.. It supplies a centre that is not in any connection.
\item Minimality of ergodic interval exchange transformations. An ergodic interval exchange transformation is minimal, so the orbit of every point is dense.. It produces an iterate of an endpoint close to the chosen centre.
\end{itemize}\medskip
\begin{corollary}\label{cor: existenceofinfiniteorbit}
Assume that $T$ is an ergodic IET. Then there exists {$\beta \in\mathcal A \setminus \{\pi_0^{-1}(1)\}$ such that $M(\beta) = +\infty$. }
\end{corollary}

\begin{corollary}
\label{cor:centers_connections}
Let $T: I \to I$ be a symmetric IET. Then any non-trivial irreducible connection of $T$ contains at most one point from the set $\left\{c_\alpha \mid \alpha \in \A \cup \{\sfrac{1}{2}\} \right\}$. In particular, there exists $\alpha \in \A$ such that $c_\alpha$ does not belong to any connection.
\end{corollary}

\begin{lemma}\label{lem: sym_endpoints}
{Let $T : I \to I$ be a symmetric IET.}
Let $\alpha\in\mathcal A\setminus\{\pi_0^{-1}(1)\}$ and $m< M(\alpha)$. Then
\begin{equation}\label{eq: connections}
\I_I\circ T(T^{-m}(\partial I_{\alpha}))={\I_I\circ (T^{-m+1}(\partial I_{\alpha}))}=T^m(\partial I_{\hat\alpha}),
\end{equation}
where $\pi_0(\hat\alpha)=\pi_0(\alpha)-1$.

In particular, the left-closed right-open interval with endpoints $T^{-m}(\partial I_{\alpha})$ and $T^m(\partial I_{\hat\alpha})$ is {$\beta$-symmetric for some $\beta\in\A$, while the left-closed right-open interval with endpoints $T^{-m+1}(\partial I_{\alpha})$ and $T^m(\partial I_{\hat\alpha})$ is $\sfrac{1}{2}$-symmetric}.

Moreover, $M(\alpha)= N(\hat\alpha)$.
\end{lemma}

\begin{lemma}
\label{lem:sym_int_disjoint}
Let $T: I \to I$ be an ergodic symmetric IET. Then there exists $\beta \in \A$ such that $c_\beta$ is not part of any non-trivial connection from $T$ and, for any $\epsilon > 0$, there exists a $\beta$-symmetric subinterval $J\subseteq I$ disjoint from the connections of $T$ satisfying $|J| < \epsilon$.
\end{lemma}

\begin{proof}
By Corollary \ref{cor:centers_connections}, there exists $\beta \in \A$ such that $c_\beta$ is not part of any connection of $T$ and, by~Corollary \ref{cor: existenceofinfiniteorbit}, there exists $\alpha \in \A \setminus \{\pi_0^{-1}(1)\}$ such that $M(\alpha) = +\infty$.

Since $T$ is ergodic and hence minimal, there exists $m \geq 1$ such that
\[
\big|T^{-m}(\partial I_\alpha) - c_\beta \big| < \sfrac{\epsilon}{2}.
\]
By taking $\epsilon$ smaller if necessary, we~may assume that $\epsilon < \min_{\delta \in \A} |I_\delta|$ and that $(c_\beta - \epsilon, c_\beta + \epsilon)$ does not contain any point from any connection.

Then, by~Lemma \ref{lem: sym_endpoints}, the left-closed right-open subinterval $J$ with endpoints $T^{-m}(\partial I_\alpha)$, $T^{m}(\partial I_{\hat \alpha})$ is $\beta$-symmetric, where $\pi_0(\hat \alpha) = \pi_0(\alpha) - 1$. In particular $J \subseteq (c_\beta - \sfrac{\epsilon}{2}, c_\beta + \sfrac{\epsilon}{2})$.
\end{proof}

\end{document}
